\section{One Request per Subgraph, per Level}
\label{sec:ch17-one-request-per-subgraph}
\index{entity fetch!batching in the router|(}%

Chapter~\ref{ch:entity-resolution-done-right} established from the outside that
\code{\_entities} is called once with a list of representations rather than once
per representation, and chapter~\ref{ch:strangling-the-monolith} watched three
of those calls happen at the same time. Neither could say what decides it, and
neither saw what the router does to the list before it sends it.

Here is the storefront query's fetch tree, taken from a trace rather than a
query plan so that the representations are visible:

\begin{minted}[fontsize=\footnotesize,breaklines]{text}
Sequence
  Single      catalog    path=''
  Parallel
    BatchEntity pricing    path='browseProducts.nodes'  reps=3
    BatchEntity inventory  path='browseProducts.nodes'  reps=3
    BatchEntity reviews    path='browseProducts.nodes'  reps=3
  BatchEntity accounts
      path='browseProducts.nodes.@.reviews.nodes.@.author'  reps=5
\end{minted}

The shape is chapter~\ref{ch:strangling-the-monolith}'s, down to that last path
with its two \code{@} steps, and needs no re-deriving. What that chapter's query
plan could not show is the column on the right, because a query plan does not
carry representations. Three of them for each of the three subgraphs
contributing to a product, which is what you would expect from three products.
And five for the authors.

\subsection*{Five representations for six authors}

\index{entity fetch!representations deduplicated}%
Three products, two reviews each, is six reviews and therefore six authors. Five
went out. One customer had written two of the six reviews, and the router
noticed.

That is the kind of observation I have been wrong about in this book before, so
rather than reason about the seed data I asked for a duplicate on purpose. Three
positions, two distinct identifiers, the first one repeated:

\begin{minted}[fontsize=\footnotesize,breaklines]{text}
asked for 3 positions holding 2 distinct ids, got 3 items back
catalog    sent 2 representations
pricing    sent 2 representations
\end{minted}

Both subgraphs were asked for two things, and the client got three, the first
two identical. The mechanism is in the loader, where each rendered
representation is hashed on its way into the outgoing buffer:

\begin{minted}[fontsize=\footnotesize,breaklines]{go}
			res.tools.keyGen.Reset()
			_, _ = res.tools.keyGen.Write(itemInput.Bytes())
			itemHash := res.tools.keyGen.Sum64()
			if existingIndex, ok := res.tools.batchHashToIndex[itemHash]; ok {
				batchStats[existingIndex] = arena.SliceAppend(res.tools.a, batchStats[existingIndex], items[i])
				continue WithNextItem
			} else {
\end{minted}

A hit on that map appends the item to the bucket belonging to the representation
already queued and jumps to the next item without writing anything. The
representation never reaches the wire a second time. Afterwards each answer is
merged onto every item in its bucket, which is how three response positions get
filled from two answers.

\index{entity fetch!two deduplications, two layers}%
Now put that next to chapter~\ref{ch:entity-resolution-done-right}. That chapter
measured four resolver calls collapsing to one key lookup, and attributed it to
GreenDonut, correctly: the DataLoader inside the subgraph deduplicates the keys
it is given. This is a different deduplication, in a different process, and it
runs first. By the time the subgraph's DataLoader sees the batch, the router has
already removed the repeats.

Both are live at once and they compose, which is worth knowing before you go
looking for a duplicate that is not arriving. A subgraph instrumented to count
distinct keys will never see the router's duplicates, so a graph where the same
customer is fetched under two names will look, from inside Accounts, like a
graph with no repetition in it at all.

\subsection*{What decides a batch, and what bounds it}

\index{entity fetch!what makes a batch}%
The condition for a batch rather than a single entity fetch is structural: the
entity jump happens at a path that is a list rather than an object. That is why
\code{browseProducts.nodes} batches and a single \code{productById} would not.
Post-processing turns that into the concrete batch form with one item template
and a comma separator, which is exactly the shape
chapter~\ref{ch:how-a-federated-query-actually-runs} printed off the wire
without knowing where it came from.

\index{entity fetch!no batch size limit}%
What bounds it is nothing I could find. The loop runs over every matching item
unconditionally, and there is no size cap in the loader. I want to be honest
about the strength of that claim: I read the loop and searched the package, and
the settings elsewhere in the router with \code{BatchSize} in the name belong to
an OpenTelemetry exporter and to the client-side request batching feature, which
are different things. What I did not do is read every path that produces the
item list, so \enquote{nothing bounds it} is the reading of one function rather
than a proof about the whole engine. If you serve a page of five hundred
entities, the \code{\_entities} call carrying five hundred representations is
what the code I read would produce.

\medskip
Which raises the obvious question of how any of this was visible at all, given
that the previous section's numbers came out of a trace. The trace is the last
subject of the chapter, and it has a sting in it.

\index{entity fetch!batching in the router|)}%
